% ============================================================================
% CHUONG 1 (BC1): BAI TOAN — UOC LUONG TAN SUAT TREN LUONG DU LIEU
% ============================================================================

\chapter{BÀI TOÁN}
\label{ch:problem}

Chương này trình bày bối cảnh thực tiễn, phát biểu toán học và phạm vi của bài toán được giải quyết trong báo cáo. Chương~\ref{ch:theory} kế tiếp sẽ đi sâu vào nền tảng lý thuyết; chương này giới hạn ở việc làm rõ \textit{cần giải cái gì, ở đâu, và dưới ràng buộc nào}.

\section{Bối cảnh và đặt vấn đề}
\label{sec:context-problem}

Bài toán mà báo cáo này giải quyết không xuất phát từ một quan sát ngẫu nhiên mà đứng trong dòng dịch chuyển kiến trúc phần mềm doanh nghiệp khoảng một thập kỷ trở lại đây: từ ứng dụng đơn khối sang vi-dịch-vụ (\textit{microservices}) đóng gói trong container và điều phối bằng Kubernetes (K8s). Sự dịch chuyển này, khởi đầu khoảng năm 2014--2015, đã trở thành chuẩn mực triển khai phần mềm doanh nghiệp; song hành với nó là một bề mặt tấn công kiểu mới mà các mô hình bảo mật chu vi truyền thống không bao phủ được --- hàng trăm endpoint nội bộ phơi bày qua mạng overlay của cluster, các lời gọi đông\nobreakdash- đông giữa dịch vụ thay vì máy khách\nobreakdash- máy chủ, vòng đời container ngắn (đo bằng phút), và sự đa dạng công nghệ trong cùng một cụm khiến bề mặt phòng thủ trở nên không đồng nhất.

Ba bài tổng quan có hệ thống (\textit{systematic literature reviews}) gần đây nhất về an toàn thông tin của hệ vi-dịch-vụ trên K8s nhất quán ở cùng một kết luận trung tâm: \textit{vẫn thiếu vắng những chuẩn bảo mật chuyên biệt và mô hình toàn diện cho kiến trúc vi-dịch-vụ}. Cụ thể, Hutasuhut và cộng sự (2024) \cite{hutasuhut2024gap} sàng lọc 487 bài báo --- chọn lại 87 bài qua phương pháp snowball --- tập trung vào các lỗ hổng phát hiện sau năm 2020, và nhấn mạnh khoảng cách kiến thức rõ rệt giữa giới học thuật và giới thực hành công nghiệp. Trước đó, Berardi và cộng sự (2022) \cite{berardi2022microservice} khảo sát 290 công bố và xác định bốn khoảng trống cốt lõi: thiếu hướng dẫn áp dụng \textit{security-by-design}, thiếu mô hình mối đe dọa phù hợp, thiếu hướng dẫn xử lý bề mặt tấn công sinh ra bởi sự đa dạng công nghệ, và các vấn đề an toàn khi di trú từ hệ đơn khối. Pereira-Vale và cộng sự (2021) \cite{pereiravale2021security} bổ sung góc nhìn đa nguồn (\textit{multivocal review}): kết hợp 36 công bố học thuật với 34 nguồn xám (báo cáo công nghiệp, blog kỹ thuật) cho thấy nhiều giải pháp công nghiệp đang được triển khai chưa có cơ sở lý thuyết được kiểm chứng, đồng thời nhiều phân tích lý thuyết chưa được hiện thực hóa trong hệ thống thực.

Đối chiếu với thực tế triển khai, hai báo cáo công nghiệp gần đây xác nhận các kết luận trên. \textit{The State of Kubernetes Security 2024} của Red Hat \cite{redhat2024k8ssec}, dựa trên khảo sát 600 chuyên gia DevOps và an toàn thông tin, ghi nhận: 67\% tổ chức đã \textit{trì hoãn hoặc làm chậm phát triển ứng dụng} do lo ngại an ninh tăng cao; 45\% gặp sự cố ở giai đoạn \textit{runtime} trong 12 tháng gần nhất; 30\% bị phạt sau sự cố và 46\% mất doanh thu hoặc khách hàng. \textit{Kubernetes Security Report 2025} của Wiz \cite{wiz2025k8ssec} bổ sung số liệu thời gian: cluster K8s mới tạo trên các dịch vụ managed (AKS, EKS) bị tấn công thử lần đầu trong vòng 18--28 phút sau khi sẵn sàng. Số liệu này cho thấy giả định ``đầu vào đối kháng'' không phải là một yêu cầu lý thuyết mà phản ánh đúng điều kiện vận hành mặc định.

Khi đối chiếu các khoảng trống nêu trên với cấu trúc một hệ giám sát mạng chạy trong cluster K8s, có thể nhận thấy một yêu cầu chung: lớp giám sát phải \textit{quan sát liên tục, ở quy mô từng gói tin, dưới ràng buộc bộ nhớ chặt}, đồng thời \textit{chịu được đầu vào đối kháng}. Một nút giám sát trung bình quan sát từ $10^5$ đến $10^7$ sự kiện mỗi giây trong các mạng sản xuất; mỗi sự kiện là một bản ghi có cấu trúc (địa chỉ IP nguồn, đích, cổng, giao thức, nhãn thời gian) và lưu lượng đi qua được mô hình hóa tự nhiên thành một \textit{luồng sự kiện}. Một thao tác cơ bản --- \textbf{ước lượng tần suất của một phần tử trên luồng} (đã quan sát bao nhiêu kết nối từ địa chỉ IP $x$? bao nhiêu yêu cầu tới cổng $p$?) --- đơn giản về mặt khái niệm nhưng trở nên thách thức vì phải đồng thời thoả ba ràng buộc vật lý: thời gian xử lý mỗi sự kiện hằng số, bộ nhớ giới hạn không tăng tuyến tính theo số phần tử phân biệt, và đảm bảo chính xác cả khi kẻ tấn công chủ động sinh nhiều khóa giả nhằm làm tràn các bảng đếm chính xác. Đây chính là miền của các thuật toán trên luồng dữ liệu (\textit{streaming algorithms}) và các cấu trúc dữ liệu xác suất (\textit{probabilistic data structures}) --- nơi đánh đổi sai số có kiểm soát lấy bộ nhớ nhỏ hơn một bậc độ lớn.

Mô hình \textit{Zero Trust} (``không tin tưởng ngầm định, luôn xác minh'') do NIST chuẩn hoá \cite{nist800207} chính là khung an ninh đặt ra yêu cầu quan sát liên tục như trên: thay vì coi mọi kết nối bên trong mạng nội bộ là tin cậy, mỗi sự kiện kết nối phải được xác thực, ủy quyền và giám sát --- không phân biệt nó đến từ bên trong hay bên ngoài. Nguyên lý này nâng mức bảo đảm an ninh bằng ba cơ chế: (i) loại bỏ các ``vùng tin cậy ngầm định'' vốn là điểm mù của tường lửa chu vi; (ii) hạn chế \textit{lateral movement} khi một nút bị xâm nhập, vì mỗi bước đi đều phải xác minh lại; (iii) thu nhỏ \textit{bán kính ảnh hưởng} của một sự cố. Việc xác minh liên tục ở quy mô từng gói tin đòi hỏi phải \textit{quan sát} được tần suất của mọi nguồn, đích, cặp (IP, cổng) trong thời gian thực --- và chính yêu cầu quan sát này sinh ra bài toán ước lượng tần suất trên luồng ở tốc độ cao và bộ nhớ hạn chế.

Báo cáo này, với phạm vi của một học phần (Công nghệ Phần mềm 1), \textit{không} trực tiếp lấp toàn bộ các khoảng trống học thuật nêu ở phần trên, mà tập trung vào một viên gạch ở tầng thấp nhất: phân tích, cài đặt và đánh giá một cấu trúc dữ liệu nền tảng (Count\nobreakdash- Min Sketch) cho một lát cắt cụ thể của bài toán giám sát --- ước lượng tần suất điểm với đảm bảo $(\varepsilon, \delta)$ và bộ nhớ hằng số. Trên nền tảng đó, Chương~\ref{ch:application} ráp lại một pipeline tinh giản \textit{tcpdump $\rightarrow$ parser $\rightarrow$ frequency estimator $\rightarrow$ alert} để chứng minh bằng thực nghiệm rằng cấu trúc dữ liệu được phân tích là đủ để phát hiện port-scan trên luồng IP thật. Việc tích hợp đầy đủ vào một hệ Zero Trust hoàn chỉnh (eBPF collector ở \textit{kernel space}, aggregator đa nút trên NATS JetStream, các dịch vụ SOC, đáp ứng tự động bằng \texttt{NetworkPolicy} qua API server Kubernetes) thuộc phạm vi của bài báo cáo ở môn tốt nghiệp Công nghệ phần mềm 2.

\section{Mô hình bài toán và yêu cầu lời giải}
\label{sec:problem-statement}

Để chuyển từ mô tả ngữ cảnh sang phát biểu hình thức, trước hết chúng ta cần thống nhất cố định năm giả thiết mô hình. \textbf{(i) Single-pass:} mỗi sự kiện được đọc đúng một lần; không có khả năng lưu trữ toàn bộ luồng để truy vấn sau. \textbf{(ii) Tốc độ đến cao:} tốc độ sự kiện vượt quá $10^5$/giây trên mỗi nút giám sát; chi phí xử lý mỗi sự kiện phải $O(1)$ hoặc $O(\log n)$. \textbf{(iii) Bộ nhớ hạn chế:} thuật toán phải chạy trong vài chục đến vài trăm KB, đủ vừa vặn với bộ nhớ cache L2 của CPU để duy trì thông lượng. \textbf{(iv) Đầu vào đối kháng:} kẻ tấn công có thể tạo ra luồng sự kiện với nhiều phần tử phân biệt (giả mạo địa chỉ IP nguồn) nhằm làm tràn các bảng băm chính xác --- thuật toán phải chịu được loại tấn công này. \textbf{(v) Chấp nhận sai số xấp xỉ:} người sử dụng chấp nhận một giá trị ước lượng $\hat{f}$ thay cho giá trị đúng $f$, miễn là sai số có thể được \textit{định lượng} và giới hạn bởi các tham số $(\varepsilon, \delta)$. Bốn giả thiết đầu là ràng buộc vật lý sinh ra từ thực tiễn vận hành; giả thiết thứ năm là một sự nới lỏng yêu cầu --- và chính sự nới lỏng này mở ra cánh cửa cho các cấu trúc sketch đạt mức bộ nhớ hằng số $O((1/\varepsilon)\log(1/\delta))$.

Trên các giả thiết đó, bài toán được phát biểu hình thức như sau. \textbf{Đầu vào} là một luồng $S = (a_1, a_2, \ldots, a_m)$ với mỗi $a_j$ thuộc một không gian phần tử $U$ có kích thước $n$. Các $a_j$ có thể trùng nhau; luồng được đọc tuần tự theo thứ tự xuất hiện, không có khả năng tua lại. \textbf{Đầu ra} là, với mỗi truy vấn $x \in U$ (có thể biết sau khi đã đọc toàn bộ luồng), một ước lượng $\hat{f}_x \in \mathbb{N}$ cho tần suất thật
\begin{equation}
f_x = \bigl|\{j : a_j = x\}\bigr|.
\label{eq:point-query-frequency}
\end{equation}

Gọi $m = \sum_x f_x$ là độ dài luồng. Thuật toán được xem là \textit{chấp nhận được} theo tiêu chí $(\varepsilon, \delta)$ nếu tồn tại hai tham số $\varepsilon, \delta \in (0, 1)$ sao cho với mọi truy vấn $x$:
\begin{equation}
\Pr\bigl[|\hat{f}_x - f_x| \le \varepsilon \cdot m\bigr] \ge 1 - \delta.
\label{eq:point-query-acceptance}
\end{equation}

Bài toán này được gọi là \textbf{Point Query} trong văn liệu thuật toán luồng và là bài toán trung tâm mà Count-Min Sketch \cite{cormode2005cms} giải quyết tối ưu. Định nghĩa hình thức chi tiết hơn (vector tần suất, moment bậc $k$, cận dưới bộ nhớ) được xây dựng ở Chương~\ref{ch:theory}; tại đây giới hạn ở mức phát biểu đủ để xác định phạm vi bài toán.

Một lời giải cho bài toán Point Query nêu trên được coi là \textbf{đạt yêu cầu} khi đồng thời thoả năm điều kiện sau. Thứ nhất, \textit{bộ nhớ giới hạn}: bộ nhớ sử dụng là $o(\min(n, m))$, lý tưởng là hằng số cố định hoặc $\mathrm{polylog}(n, m)$. Thứ hai, \textit{thời gian cập nhật hằng số}: thao tác xử lý mỗi sự kiện là $O(1)$ hoặc tối đa $O(\log(1/\delta))$, không phụ thuộc $m$ hay $n$. Thứ ba, \textit{đảm bảo có chứng minh}: sai số phải có cận lý thuyết được chứng minh từ các giả thiết về hàm băm, không dựa vào phân bố dữ liệu cụ thể. Thứ tư, \textit{xác nhận thực nghiệm}: cài đặt phải được kiểm thử đơn vị cho từng tính chất cốt lõi, benchmark cho thông lượng và bộ nhớ, và trace từng bước để kiểm chứng hoạt động. Thứ năm, \textit{có thể tái lập}: mọi tham số, dữ liệu đầu vào, seed ngẫu nhiên phải được ghi lại để kết quả báo cáo có thể được tái tạo từ mã nguồn. Năm điều kiện này sẽ là khung tham chiếu để đánh giá và so sánh các phương án ở Chương~\ref{ch:compare}.

\section{Khảo sát công trình liên quan}
\label{sec:related-work}

Các công trình nền tảng cho bài toán Point Query nói riêng và bài toán Heavy Hitters nói chung được nhóm thành hai mạch: các kết quả giới hạn lý thuyết cùng với thuật toán deterministic, và các cấu trúc sketch xác suất cùng với khảo sát thực nghiệm đi kèm.

Ở mạch thứ nhất, Alon, Matias và Szegedy (1996) \cite{alon1996space} thiết lập giới hạn dưới cho không gian các thuật toán ước lượng moment tần suất trên luồng: $\Omega(1/\varepsilon)$ cho thuật toán deterministic và $\Omega((1/\varepsilon)\log(1/\delta))$ cho thuật toán randomized. Công trình này (nhận Giải thưởng Gödel 2005) định hình toàn bộ miền và khẳng định mọi lời giải sau đó --- kể cả Count-Min Sketch --- không thể vượt qua ngưỡng này. Đặt trong khung đó, Misra và Gries (1982) \cite{misra1982finding} là thuật toán deterministic cổ điển: với bộ nhớ $O(k)$ giữ lại tối đa $k$ phần tử ``nặng'' trong luồng, thuật toán đạt được cận dưới $\Omega(1/\varepsilon)$ nhưng chỉ cho ước lượng \textit{dưới mức} và giải bài toán Heavy Hitters chứ không phải Point Query tổng quát --- khiến nó trở thành một \textit{điểm so sánh} hơn là lời giải trực tiếp cho bài toán báo cáo đặt ra.

Ở mạch thứ hai, Cormode và Muthukrishnan (2005) \cite{cormode2005cms} giới thiệu Count-Min Sketch --- cấu trúc sketch hai chiều sử dụng $d$ hàm băm pairwise-independent. Với bộ nhớ $O((1/\varepsilon)\log(1/\delta))$, CMS đạt đảm bảo $(\varepsilon, \delta)$ cho Point Query và đạt tối ưu so với cận dưới ở mạch thứ nhất. Biến thể \textit{Conservative Update} do Estan và Varghese \cite{estan2003conservative} đề xuất trong bối cảnh đo lưu lượng mạng giảm đáng kể sai số thực tế trên phân phối Zipfian --- đặc trưng của lưu lượng Internet. Để so sánh đối chiếu, Charikar, Chen và Farach-Colton (2004) \cite{charikar2004count} đề xuất Count Sketch: một biến thể cho ước lượng \textit{không thiên lệch} bằng cách sử dụng hàm sign $s_i$ và lấy median thay vì min, đánh đổi bằng bộ nhớ $O((1/\varepsilon^2)\log(1/\delta))$ --- lớn hơn CMS một bậc $1/\varepsilon$ --- nên phù hợp với các toán tử đại số (inner product, join size) hơn là Point Query. Cuối cùng, Cormode và Hadjieleftheriou (2008) \cite{cormode2008finding} thực hiện so sánh thực nghiệm toàn diện giữa các thuật toán tần suất phổ biến trên nhiều tập dữ liệu khác nhau, cung cấp thông tin định lượng về thông lượng và độ chính xác trong các kịch bản thực tế --- là cơ sở để so sánh cài đặt trong chương đánh giá.

\section{Phạm vi và định hướng giải quyết}
\label{sec:scope-and-choice}

Báo cáo giải quyết bài toán \textbf{Point Query} trên mô hình luồng \textit{cash-register} (insert-only), với mục tiêu chính là cài đặt và đánh giá Count-Min Sketch và các thuật toán thay thế (Misra--Gries, Hash Map chính xác). Các vấn đề nằm ngoài phạm vi gồm: mô hình \textit{turnstile} cho phép giảm counter (do các thuật toán như Count Sketch hoặc CMS có sign hỗ trợ, không cần thiết khi đầu vào chỉ có chèn); bài toán Top-$k$ / Heavy Hitters đầy đủ liệt kê được các phần tử có tần suất cao nhất (Misra--Gries giải một phần, CMS cần kết hợp cấu trúc phụ); các bài toán Range Query, Inner Product Estimation và Join Size Estimation thuộc lớp toán tử đại số trên sketch; và các thuật toán phụ khác trong hệ thống Zero-Trust (Tarjan SCC, BFS k-hop, LPM Trie) --- được đề cập ở phần hướng phát triển của Chương~\ref{ch:eval} và là nội dung của đồ án tốt nghiệp.

Để người đọc có một định hướng rõ ràng trước khi bước vào phần lý thuyết, báo cáo công bố trước kết luận: \textbf{Count-Min Sketch (CMS) với biến thể Conservative Update} là lời giải được chọn cho bài toán Point Query nêu trên. Các chương sau đi theo lộ trình tuyến tính dẫn tới kết luận này. Chương~\ref{ch:theory} xây dựng nền tảng lý thuyết của CMS: cấu trúc ma trận $d \times w$, họ hàm băm pairwise-independent, định lý Cormode--Muthukrishnan và chứng minh cận sai số $(\varepsilon, \delta)$. Chương~\ref{ch:compare} so sánh CMS với ba phương án thay thế (Hash Map chính xác, Misra--Gries, Count Sketch) trên bốn yêu cầu cốt lõi --- bộ nhớ hằng số, cập nhật $O(1)$, sai số một phía có thể định lượng và khả năng hợp nhất song song --- rồi chỉ ra lý do cụ thể để loại từng phương án. Chương~\ref{ch:impl} trình bày cài đặt CMS trong thư viện \texttt{libdsa} và đo thông lượng, sai số thực tế trên dữ liệu benchmark. Chương~\ref{ch:eval} tổng kết kết quả và xác nhận lại lựa chọn bằng số liệu đo được. Cách trình bày này cho phép người đọc đối chiếu từng đặc tính của CMS với yêu cầu bài toán ngay khi gặp, thay vì phải chờ tới kết luận cuối để biết giải thuật nào được chọn và vì sao.

\section{Kết luận chương}

Chương này đã: (i) đặt bài toán giám sát mạng dưới khung Zero Trust trong bối cảnh các nghiên cứu khoảng trống trên Kubernetes, và chỉ ra rằng việc xác minh liên tục ở quy mô gói tin sinh ra một bài toán tính toán cụ thể --- ước lượng tần suất trên luồng dưới ràng buộc bộ nhớ và đầu vào đối kháng; (ii) phát biểu hình thức bài toán \textbf{Point Query} trên mô hình cash-register với tiêu chí $(\varepsilon, \delta)$ và năm điều kiện chấp nhận lời giải; (iii) khảo sát các công trình nền tảng (Alon--Matias--Szegedy, Misra--Gries, Cormode--Muthukrishnan, Estan--Varghese, Charikar--Chen--Farach-Colton, Cormode--Hadjieleftheriou) làm cơ sở cho phần so sánh ở chương sau; (iv) công bố trước kết luận: \textbf{Count-Min Sketch với Conservative Update là phương án được chọn}, đồng thời chỉ rõ phạm vi (Point Query, cash-register) và các vấn đề nằm ngoài phạm vi báo cáo. Bốn nội dung này thiết lập đầy đủ ngữ cảnh và yêu cầu; các chương kế tiếp lần lượt xây dựng nền tảng lý thuyết, so sánh phương án, cài đặt và đánh giá để chứng minh lựa chọn vừa nêu.
